\begin{table}[t]
\centering
\caption{Detection ablation: only the peak detector is swapped; area
integration and \rf{}'s assignment stage are held identical, so scores
reflect detection alone. Columns: mean and worst-case composition L1
(at.\%) over the 14 datasets every detector completed, count of
catastrophic failures (L1$>$25\,at.\%), mean fabricated signal, and on
the synthetic specimens peak precision and median mass error. \rf{}'s
multi-resolution detector is the only one that never fails
catastrophically; its mass accuracy reflects raw-event centroiding and
family recentering that bin-apex detectors cannot match. Bold = best.}
\label{tab:detection_ablation}
\small
\begin{tabular}{lrrrrrr}
\toprule
Detector & mean L1 & worst L1 & \#L1$>$25 & fab. & Prec. & $|\Delta m|$ ppm \\
\midrule
\rf{} (multi-resolution) & \textbf{10.36} & \textbf{53.80} & 2 & \textbf{1.84} & \textbf{0.984} & \textbf{2.2} \\
Naive prominence & 28.51 & 184.93 & 4 & 10.63 & 0.959 & 113.9 \\
Savitzky--Golay prominence & 22.97 & 87.26 & 5 & 6.84 & 0.806 & 139.0 \\
CWT ridge lines & 30.90 & 184.93 & 4 & 11.99 & 0.963 & 113.9 \\
Persistent homology & 33.57 & 184.93 & 4 & 13.25 & 0.935 & 113.9 \\
Derivative zero-crossing & 23.76 & 81.12 & 5 & 7.64 & 0.867 & 139.0 \\
\bottomrule
\end{tabular}
\end{table}
